\documentclass[a4paper,11pt,reqno]{amsart}

\usepackage[T1]{fontenc}
\usepackage{lmodern}
\usepackage{microtype}
\usepackage{amsmath,amssymb,amsthm}
\usepackage{xcolor}
\usepackage[a4paper,left=21mm,right=21mm,top=17mm,bottom=19mm]{geometry}

\newtheorem{fact}{Fact}
\newtheorem*{mainlemma}{Main Lemma}
\newtheorem*{maintheorem}{Main Theorem}

\newcommand{\Z}{\mathbb Z}
\newcommand{\U}{\operatorname U}
\newcommand{\w}{\varpi}
\newcommand{\dlt}{\delta}
\newcommand{\g}{\gamma}

\setlength{\abovedisplayskip}{5pt plus 2pt minus 2pt}
\setlength{\belowdisplayskip}{5pt plus 2pt minus 2pt}
\setlength{\abovedisplayshortskip}{3pt plus 2pt}
\setlength{\belowdisplayshortskip}{3pt plus 2pt minus 2pt}

\title{Centrality of Lie Dimension Subrings}
\author{}
\date{}

\begin{document}

\maketitle
\vspace{-1.2em}

Let \(L\) be a Lie algebra over \(\Z\). Let \(\U(L)\) be its universal
 enveloping algebra, let \(\iota\colon L\to\U(L)\) be the canonical Lie
homomorphism, and let \(\varepsilon\colon\U(L)\to\Z\) be the augmentation.
Put \(\w(L)=\ker\varepsilon\), and, for \(n\geq1\), define
\[
  \dlt_n(L)=\iota^{-1}\bigl(\w(L)^n\bigr),\qquad
  \g_1(L)=L,\qquad \g_{n+1}(L)=[\g_n(L),L].
\]
For additive subgroups \(A,B\subseteq L\), let \([A,B]\) denote the
additive subgroup generated by the elements \([a,b]\), where \(a\in A\)
and \(b\in B\). All iterated commutators are left-normed.

\begin{fact}\label{fact:lcs}
For every \(m\geq1\), the subgroup \(\g_m(L)\) is an ideal,
\(\g_{m+1}(L)\subseteq\g_m(L)\), and \(\g_m(L)\) is generated by the
commutators \([x_1,\ldots,x_m]\) with \(x_i\in L\).
\end{fact}

\begin{fact}\label{fact:augmentation}
For every \(n\geq1\), the additive group \(\w(L)^n\) is generated by
\(\iota(x_1)\cdots\iota(x_r)\), where \(r\geq n\) and \(x_i\in L\).
\end{fact}

\begin{fact}\label{fact:standard}
For every \(n\geq1\), one has \(\g_n(L)\subseteq\dlt_n(L)\).
\end{fact}

\begin{mainlemma}
There is a right \(\U(L)\)-module structure on \(L\) satisfying
\[
  z\cdot\iota(x)=[z,x]\qquad(z,x\in L).
\]
For this action,
\[
  L\cdot\w(L)^n=\g_{n+1}(L)\qquad(n\geq1),
\]
where \(L\cdot\w(L)^n\) is the subgroup generated by all \(z\cdot u\) with
\(z\in L\) and \(u\in\w(L)^n\).
\end{mainlemma}

\begin{proof}
Let the tensor algebra of the underlying abelian group of \(L\) act on
\(L\), extending \(\Z\)-linearly the rules
\[
  z\cdot1=z,\qquad z\cdot(x_1\cdots x_r)=[z,x_1,\ldots,x_r].
\]
For \(x,y,z\in L\), the Jacobi identity gives
\[
 z\cdot\bigl(xy-yx-[x,y]\bigr)
 =[[z,x],y]-[[z,y],x]-[z,[x,y]]=0.
\]
Thus the defining ideal of \(\U(L)\) acts trivially, so the action factors
through \(\U(L)\).

If \(r\geq n\), then Fact~\ref{fact:lcs} gives
\[
 z\cdot\iota(x_1)\cdots\iota(x_r)
 =[z,x_1,\ldots,x_r]\in\g_{r+1}(L)\subseteq\g_{n+1}(L).
\]
By Fact~\ref{fact:augmentation}, this proves
\(L\cdot\w(L)^n\subseteq\g_{n+1}(L)\). Conversely, Fact~\ref{fact:lcs}
shows that \(\g_{n+1}(L)\) is generated by
\[
 [x_0,x_1,\ldots,x_n]
 =x_0\cdot\iota(x_1)\cdots\iota(x_n)\in L\cdot\w(L)^n.
\]
\end{proof}

\begin{center}
\setlength{\fboxsep}{7pt}
\fcolorbox{green!45!black}{green!9}{%
\begin{minipage}{0.91\linewidth}
\begin{maintheorem}
For every Lie algebra \(L\) over \(\Z\) and every \(n\geq1\),
\[
  [\dlt_n(L),L]=\g_{n+1}(L).
\]
\end{maintheorem}
\end{minipage}}
\end{center}
\vspace{-0.4em}

\begin{proof}
Let \(a\in\dlt_n(L)\) and \(x\in L\). Since
\(\iota(a)\in\w(L)^n\), the Main Lemma gives
\[
 [x,a]=x\cdot\iota(a)\in\g_{n+1}(L).
\]
By antisymmetry, \([a,x]\in\g_{n+1}(L)\), and hence
\([\dlt_n(L),L]\subseteq\g_{n+1}(L)\). Conversely, Fact~\ref{fact:standard}
gives
\[
 \g_{n+1}(L)=[\g_n(L),L]\subseteq[\dlt_n(L),L].
\]
\end{proof}

\end{document}
